\section{Foundations (c.~1950--1979)}
\label{sec:foundations}

This section separates what the distributed-memory compilers of the late 1980s inherited from what they invented. Nothing here is a distributed-memory compiler. The machines of the period had a single address space or lock-step SIMD control, and placing data in memories was not a compiler concern. Every foundations paper discussed here has a card in the repository, and its content claims cite that card.

\subsection{Automatic programming}
\fact{The FORTRAN project, begun in 1954, aimed to reduce the effort of preparing scientific problems for the IBM~704, and its translator turned algebraic statements into 704 programs~\cite{backus1957fortran}.} \fact{The authors report one job: four hours to write 47 statements, six minutes to compile them into about 1000 instructions, and one correction. They report no controlled benchmark and no run times~\cite{backus1957fortran}.} \ours{We take this as the first instance of the bargain that every later system renews: the programmer writes a higher-level program, and the compiler makes the low-level decisions.}

\subsection{Dependence analysis and loop parallelism}
\fact{Lamport asks how the iterations of a DO nest can run concurrently and still give the sequential results. His coordinate and hyperplane methods are analytic: the paper gives step counts for its loop examples and no measurements~\cite{lamport1974parallel}.} \fact{Kuck and colleagues propose dependence graphs, with loop, flow, output, antidependence and input arcs, as the vehicle for compiler optimisations. Their only quantitative transformation results are page reductions from loop distribution, and no speedups are reported~\cite{kuck1981dependence}.} \fact{Allen and Kennedy's translator rewrites FORTRAN DO loops as array statements, with dependence as the criterion for a correct translation. Its only measurement is translation time, about one minute of CPU time for PFC against about ten for PARAFRASE~\cite{allen1987automatic}.} \infer{Distributed-memory compilers reused this dependence representation without change for the iteration and access analysis; what they added was an ownership function. Section~\ref{sec:mapping} shows this for the Fortran~D and PARADIGM lineage.}

\subsection{Machine abstractions}
\fact{The ILLIAC~IV is a parallel-array computer of 256 processing elements, described for large matrix and grid problems~\cite{barnes1968illiac}, and the CRAY-1 is a vector computer whose compiler and operating system the paper describes~\cite{russell1978cray1}.} \infer{These two machine classes, array and vector, are the regular-parallel targets that the early loop compilers addressed.} \fact{The 1977 MIT workshop report on data-flow computers and program organisation collects talks on programming data-flow machines. It reports simulated processor allocation and waiting times and speedups for unoptimised graphs, but no controlled evaluation~\cite{misunas1977workshop}.} \fact{Kung argues that special-purpose systems are often limited by I/O bandwidth rather than cell speed, and proposes systolic arrays, pipelines of simple cells that reuse each fetched item~\cite{kung1982systolic}. The paper reports no experiments; its evidence is analytic, with I/O-rate and utilisation calculations for named designs~\cite{kung1982systolic}.} \fact{The Cosmic Cube gave each node its own storage and has its processes communicate by message passing, not shared variables~\cite{seitz1985cosmic}.} \infer{Systolic arrays and the Cosmic Cube are the architectural ancestors of the spatial and distributed targets discussed in Sections~\ref{sec:modern} and~\ref{sec:history}.}

\subsection{Data-parallel programming}
\fact{Hillis and Steele describe data-parallel algorithms with one virtual processor per element, typically running in logarithmic time. Their notation is a parallel pseudo-code over virtual processors rather than a language~\cite{hillis1986data}.} \fact{Loveman et al.\ describe HPF as Fortran~90 with mapping directives and a few first-class constructs such as FORALL, with array operations as its base~\cite{loveman1993hpf}.} \infer{The whole-array style of Fortran~90 is the lineage from the data-parallel model to HPF. The link is our inference; the HPF specification itself was not read.}

\subsection{Synthesis}
\infer{The foundations supplied three things to the next era: a trusted pipeline for transforming loops, a legality test for reordering them, and machine models with regular parallel structure. What they did not supply was any representation of where data resides. That omission became the central compiler problem of the late 1980s, once machines with private memories appeared~\cite{seitz1985cosmic}.}
